\section{SPECJALNE MOŻLIWOŚCI NIEKTÓRYCH POTWORÓW}

Ogólna zasada:

Niektóre potwory posiadają specjalne możliwości, które są wykorzystywane podczas walki z oddziałem.

Zasady szczegółowe:

[15.1] Kronki śmierdzą w sposób tak straszny, iż może to przyprawić o chorobę.
Jeżeli oddział napotka kronki i rozpoczyna z nimi walkę każdy z jego członków musi podjąć próbę odparcia czaru.
Postać, która tej próby nie przejdzie w sposób pomyślny zaczyna chorować i jej zręczność jest obniżona w czasie tej walki o 2 punkty.
Po walce z kronkami postać zdrowieje.

[15.2] Chimera, atakując postać w normalny sposób, może dodatkowo zionąć ogniem.
Zaatakowana postać musi próbować odeprzeć czar, a jeżeli próba będzie nieudana otrzymuje, poza ranami, odniesionymi na skutek normalnego ataku chimery, jedną dodatkową ranę od ognia.

[15.3] Meduza może zamienić zaatakowaną postać w kamień.
Przy każdym spotkaniu z Meduzą należy rzucać 1K6.
Jeżeli zostanie wyrzucone 6 oczek zaatakowana przez Meduzę postać zamienia się w kamień.
Żeton, reprezentujący tę postać powinien zostać usunięty z Szyku oddziału.

Jeżeli któryś z członków oddziału zna czar ,,Ożywienie kamienia'' może, bezpośrednio po zakończeniu walki, przywrócić zamienioną w kamień postać do życia.
Jeżeli nikt nie zna tego czaru postać ta jest uważana za martwą.

[15.4] Troll posiada zdolność regeneracji.
Po zakończeniu każdej trzeciej fazy ,,atak potwora'' Troll zmniejsza o jeden liczbę ran, które odniósł.

[15.5] Wampir posiada zdolność rzucania ,,za darmo'' czaru ,,urok'' ale tylko podczas dwóch pierwszych rund walki.
Jeżeli zaatakowanej tym czarem postaci nie uda się go odeprzeć dostaje się ona we władzę Wampira i jest umieszczana w pierwszym szeregu szyku potworów.
Czar może być złamany tylko przez śmierć Wampira lub przez czar ,,uwolnienie''.
Postać, znajdująca się pod wpływem uroku nie może rzucać czarów, jednak musi w inny sposób atakować oddział.

[15.6] Każda rana zadana Hydrze zwiększa jej zręczność o jeden punkt.
Tak więc w toku walki możliwości bojowe Hydry rosną wprost proporcjonalnie do liczby odniesionych ran.
Hydra się uśmiercana w normalny sposób, tzn. jeżeli liczba zadanych jej ran zrówna się ze współczynnikiem siła.

[15.7] Czarnoksiężnicy (wraz z Bezimiennym) znają jeden czar~-- ,,błyskawica''.
Używają go w trakcie każdej fazy ,,atak potwora''.
Czarnoksiężnik nie może użyć tego czaru, jeżeli związany z nim koszt spowodowałby jego śmierć.
